% Mathématiques 6e — cours V3
% Source LaTeX rendue côté Astro/KaTeX.

\section{Objectifs}

Dans ce chapitre, on apprend à :

\begin{itemize}
\item comprendre une fraction comme \textbf{nombre}, \textbf{quotient} et \textbf{opérateur multiplicatif} ;
\item relier l’écriture fractionnaire à une division exacte ;
\item utiliser une fraction pour résoudre une égalité multiplicative à trou ;
\item placer, repérer, comparer, encadrer et ordonner des fractions ;
\item reconnaître et produire des fractions égales ;
\item additionner et soustraire des fractions dans les cas étudiés ;
\item multiplier une fraction par un nombre entier ;
\item résoudre des problèmes avec des fractions ;
\item comprendre une proportion et l’exprimer sous forme de pourcentage ;
\item appliquer un pourcentage à une quantité ou à un nombre.
\end{itemize}

\section{1. Une fraction représente un nombre}

Une fraction s’écrit sous la forme :

\[
\dfrac{a}{b}
\]

où $a$ est un entier et $b$ un entier non nul.

Le nombre $a$ est appelé \textbf{numérateur}.

Le nombre $b$ est appelé \textbf{dénominateur}.

Le dénominateur indique en combien de parts égales on partage l’unité de référence ; le numérateur indique combien de ces parts sont considérées.

Par exemple :

\[
\dfrac{3}{4}
\]

peut représenter trois parts lorsque l’unité a été partagée en quatre parts égales.

Mais en 6e, on va plus loin : une fraction ne désigne pas seulement une « partie d’un tout ». Elle est un \textbf{nombre à part entière}.

\section{2. Le nouveau sens de quotient}

La fraction

\[
\dfrac{a}{b}
\]

représente le résultat exact de la division de $a$ par $b$.

On écrit :

\[
\boxed{\dfrac{a}{b}=a\div b}
\]

avec $b\neq0$.

Ainsi :

\[
\dfrac{3}{4}=3\div4.
\]

Cette égalité est essentielle : la barre de fraction peut être lue comme un symbole de division.

\subsection{Exemple}

\[
\dfrac{7}{2}=7\div2=3{,}5.
\]

Donc la fraction $\dfrac{7}{2}$ représente le nombre décimal $3{,}5$.

En revanche, certaines divisions ne donnent pas une écriture décimale finie.

Par exemple :

\[
\dfrac{2}{3}=2\div3.
\]

La fraction $\dfrac{2}{3}$ est alors une écriture exacte du quotient, même lorsque l’écriture décimale n’est pas finie.

\section{3. Comprendre le quotient par un partage}

Considérons un partage où $3$ unités identiques doivent être réparties équitablement entre $4$ personnes.

\coursefigure{/images/mathematiques/6eme/fractions-partage-quotient.svg}{Schéma de trois bandes identiques partagées chacune en quatre parts égales pour illustrer une fraction comme quotient.}{Répartir $3$ unités en $4$ parts égales conduit au quotient $\dfrac{3}{4}$.}

Chaque personne reçoit donc :

\[
\dfrac{3}{4}
\]

d’une unité.

Cette représentation permet de comprendre :

\[
3\div4=\dfrac{3}{4}.
\]

La fraction est donc bien le \textbf{résultat exact d’un partage}.

\section{4. Définition du quotient}

Le quotient de $a$ par $b$, avec $b\neq0$, est le nombre qui, multiplié par $b$, donne $a$.

Autrement dit :

\[
\boxed{
b\times\dfrac{a}{b}=a
}
\]

et, par commutativité de la multiplication :

\[
\boxed{
\dfrac{a}{b}\times b=a
}.
\]

Par exemple :

\[
4\times\dfrac{3}{4}=3.
\]

La fraction $\dfrac{3}{4}$ est donc exactement le nombre qui complète l’égalité :

\[
4\times\boxed{\phantom{\dfrac{3}{4}}}=3.
\]

\section{5. Égalités multiplicatives à trou}

Le sens quotient permet de résoudre des égalités du type :

\[
b\times x=a.
\]

Le nombre cherché est :

\[
x=\dfrac{a}{b}.
\]

\subsection{Exemple développé 1 — trouver le nombre manquant}

On cherche $x$ tel que :

\[
5\times x=3.
\]

Le nombre qui, multiplié par $5$, donne $3$ est :

\[
x=\dfrac{3}{5}.
\]

Vérification :

\[
5\times\dfrac{3}{5}=3.
\]

Dans ce cas, une écriture fractionnaire est particulièrement naturelle.

\section{6. Une fraction peut être inférieure, égale ou supérieure à $1$}

Une fraction n’est pas nécessairement inférieure à $1$.

Si le numérateur est inférieur au dénominateur :

\[
\dfrac{3}{5}<1.
\]

Si le numérateur est égal au dénominateur :

\[
\dfrac{5}{5}=1.
\]

Si le numérateur est supérieur au dénominateur :

\[
\dfrac{7}{5}>1.
\]

On peut écrire :

\[
\dfrac{7}{5}=1+\dfrac{2}{5}.
\]

Cette écriture est appelée ici \textbf{nombre mixte} : une partie entière à laquelle on ajoute une fraction inférieure à $1$.

\section{7. Fraction entière, décimale ou non décimale}

Une fraction peut représenter différents types de nombres.

\subsection{Un nombre entier}

\[
\dfrac{12}{4}=3.
\]

\subsection{Un nombre décimal non entier}

\[
\dfrac{3}{4}=0{,}75.
\]

\subsection{Un nombre qui n’a pas d’écriture décimale finie}

\[
\dfrac{2}{3}
\]

ne possède pas d’écriture décimale finie.

Il est donc important de ne pas croire qu’une fraction est toujours « un nombre à virgule caché ».

L’écriture fractionnaire peut être l’écriture exacte la plus adaptée.

\section{8. Placer une fraction sur une demi-droite graduée}

Pour placer une fraction comme :

\[
\dfrac{5}{4},
\]

il faut que la graduation permette de partager chaque unité en quatre parts égales.

\coursefigure{/images/mathematiques/6eme/fractions-droite-graduee.svg}{Demi-droite graduée dont chaque unité est partagée en quatre parts égales afin de placer plusieurs fractions.}{Lorsque chaque unité est partagée en $4$ parts égales, chaque petite graduation correspond à $\dfrac{1}{4}$.}

On a :

\[
\dfrac{5}{4}
=
1+\dfrac{1}{4}.
\]

Le point d’abscisse $\dfrac{5}{4}$ se situe donc un quart après $1$.

\section{9. Graduer un segment}

Graduer un segment signifie choisir puis construire des intervalles égaux.

Si une unité doit être partagée en $5$ parts égales, chaque petite graduation vaut :

\[
\dfrac{1}{5}.
\]

Les positions successives sont alors :

\[
0,\quad
\dfrac{1}{5},\quad
\dfrac{2}{5},\quad
\dfrac{3}{5},\quad
\dfrac{4}{5},\quad
1.
\]

Pour un dénominateur $b$, une unité est partagée en $b$ intervalles égaux.

\section{10. Fractions égales}

Deux fractions différentes peuvent représenter le même nombre.

Par exemple :

\[
\dfrac{1}{2}
=
\dfrac{2}{4}
=
\dfrac{3}{6}.
\]

On obtient une fraction égale en multipliant le numérateur et le dénominateur par un même entier non nul.

Si $k\neq0$ :

\[
\boxed{
\dfrac{a}{b}
=
\dfrac{a\times k}{b\times k}
}.
\]

Par exemple :

\[
\dfrac{3}{5}
=
\dfrac{3\times4}{5\times4}
=
\dfrac{12}{20}.
\]

On peut également reconnaître une fraction égale lorsque le numérateur et le dénominateur ont été divisés par un même facteur.

\section{11. Comparer des fractions de même dénominateur}

Lorsque deux fractions ont le même dénominateur, on compare leurs numérateurs.

Par exemple :

\[
\dfrac{5}{8}<\dfrac{7}{8}.
\]

Les parts ont la même taille ; on compare seulement le nombre de parts.

\section{12. Comparer des fractions de même numérateur}

Lorsque deux fractions positives ont le même numérateur, la fraction ayant le \textbf{plus petit dénominateur} est la plus grande.

Par exemple :

\[
\dfrac{3}{4}>\dfrac{3}{7}.
\]

En effet, partager une unité en $4$ parts donne des parts plus grandes que la partager en $7$ parts.

Cette comparaison doit être comprise par le sens de la fraction, pas mémorisée isolément.

\section{13. Comparer grâce à des fractions égales}

Pour comparer :

\[
\dfrac{3}{4}
\qquad\text{et}\qquad
\dfrac{7}{8},
\]

on peut écrire :

\[
\dfrac{3}{4}
=
\dfrac{6}{8}.
\]

Ainsi :

\[
\dfrac{6}{8}<\dfrac{7}{8},
\]

donc :

\[
\boxed{
\dfrac{3}{4}<\dfrac{7}{8}
}.
\]

En 6e, on privilégie des transformations simples et compréhensibles.

\section{14. Encadrer une fraction}

Encadrer une fraction consiste à trouver deux nombres entre lesquels elle se situe.

Par exemple :

\[
1<\dfrac{7}{4}<2.
\]

car :

\[
\dfrac{7}{4}
=
1+\dfrac{3}{4}.
\]

De même :

\[
2<\dfrac{11}{5}<3.
\]

Un encadrement permet de situer rapidement l’ordre de grandeur d’une fraction.

\section{15. Ordonner des fractions et des nombres mixtes}

Pour ordonner plusieurs nombres, on peut :

\begin{enumerate}
\item repérer leur partie entière ;
\item utiliser une demi-droite graduée ;
\item transformer certaines fractions en fractions égales ;
\item utiliser des valeurs de référence comme $0$, $\dfrac{1}{2}$, $1$ ou $2$.
\end{enumerate}

\subsection{Exemple}

Comparons :

\[
\dfrac{3}{4},
\qquad
1+\dfrac{1}{4},
\qquad
\dfrac{5}{4},
\qquad
\dfrac{7}{8}.
\]

On remarque :

\[
1+\dfrac{1}{4}
=
\dfrac{5}{4}.
\]

Et :

\[
\dfrac{3}{4}
=
\dfrac{6}{8}.
\]

Donc :

\[
\boxed{
\dfrac{3}{4}
<
\dfrac{7}{8}
<
\dfrac{5}{4}
=
1+\dfrac{1}{4}
}.
\]

\section{16. Additionner des fractions}

Lorsque deux fractions ont le même dénominateur :

\[
\dfrac{a}{b}+\dfrac{c}{b}
=
\dfrac{a+c}{b}.
\]

Par exemple :

\[
\dfrac{2}{7}+\dfrac{3}{7}
=
\dfrac{5}{7}.
\]

On additionne les numérateurs car les parts ont la même taille.

Il serait faux d’écrire :

\[
\dfrac{2}{7}+\dfrac{3}{7}
=
\dfrac{5}{14}.
\]

Le dénominateur reste $7$.

\section{17. Additionner après avoir créé un dénominateur commun simple}

Dans certains cas simples :

\[
\dfrac{1}{2}+\dfrac{1}{4}.
\]

On utilise :

\[
\dfrac{1}{2}=\dfrac{2}{4}.
\]

Donc :

\[
\dfrac{1}{2}+\dfrac{1}{4}
=
\dfrac{2}{4}+\dfrac{1}{4}
=
\dfrac{3}{4}.
\]

Le principe est de travailler avec des parts de même taille.

\section{18. Soustraire des fractions}

Lorsque les dénominateurs sont identiques :

\[
\dfrac{a}{b}-\dfrac{c}{b}
=
\dfrac{a-c}{b}.
\]

Par exemple :

\[
\dfrac{7}{10}-\dfrac{3}{10}
=
\dfrac{4}{10}.
\]

On peut ensuite reconnaître une fraction égale plus simple :

\[
\dfrac{4}{10}=\dfrac{2}{5}.
\]

\section{19. Multiplier une fraction par un entier}

Le produit d’une fraction par un entier peut s’interpréter comme une addition répétée.

Par exemple :

\[
3\times\dfrac{2}{5}
=
\dfrac{2}{5}+\dfrac{2}{5}+\dfrac{2}{5}
=
\dfrac{6}{5}.
\]

On peut donc écrire :

\[
\boxed{
n\times\dfrac{a}{b}
=
\dfrac{n\times a}{b}
}.
\]

Grâce à la commutativité :

\[
\dfrac{a}{b}\times n
=
n\times\dfrac{a}{b}.
\]

\section{20. La fraction comme opérateur multiplicatif}

Prendre une fraction d’un nombre revient à effectuer une multiplication.

Par exemple, prendre :

\[
\dfrac{3}{4}
\]

de $20$ revient à calculer :

\[
\dfrac{3}{4}\times20.
\]

On peut raisonner en deux étapes :

\[
20\div4=5,
\]

puis :

\[
5\times3=15.
\]

Donc :

\[
\boxed{
\dfrac{3}{4}\times20=15
}.
\]

Le dénominateur indique le partage, puis le numérateur indique combien de parts sont prises.

\section{21. Exemple développé 2 — fraction d’une grandeur}

Une course mesure $18\ \text{km}$. Un participant a parcouru :

\[
\dfrac{2}{3}
\]

de la distance.

On calcule :

\[
\dfrac{2}{3}\times18.
\]

Un tiers de $18$ vaut :

\[
18\div3=6.
\]

Deux tiers valent alors :

\[
2\times6=12.
\]

Le participant a donc parcouru :

\[
\boxed{12\ \text{km}}.
\]

\section{22. Résoudre un problème avec des fractions}

Lorsqu’une fraction intervient dans un problème, il faut identifier son rôle.

Elle peut représenter :

\begin{itemize}
\item une part d’un tout ;
\item un quotient exact ;
\item une position sur une demi-droite ;
\item une fraction d’une quantité ;
\item une proportion.
\end{itemize}

\subsection{Méthode}

\begin{enumerate}
\item Identifier \textbf{ce que représente l’unité}.
\item Identifier le rôle du dénominateur.
\item Identifier le rôle du numérateur.
\item Choisir une représentation si nécessaire.
\item Effectuer le calcul.
\item Interpréter le résultat dans le contexte.
\end{enumerate}

\section{23. Inventer un problème correspondant à une fraction}

Savoir inventer un problème permet de vérifier qu’on comprend le sens d’une écriture.

Par exemple, l’expression :

\[
\dfrac{3}{5}\times40
\]

peut correspondre à la situation suivante :

\begin{remarque}
Une bibliothèque possède $40$ livres sur une étagère. Les $\dfrac{3}{5}$ sont des romans. Combien de romans y a-t-il ?
\end{remarque}

Le calcul donne :

\[
\dfrac{3}{5}\times40
=
24.
\]

Un bon problème doit donner un sens clair à l’unité, à la fraction et à la quantité finale.

\section{24. Comprendre un pourcentage}

Un pourcentage exprime une proportion \textbf{pour $100$}.

Par exemple :

\[
25\,\%
=
\dfrac{25}{100}.
\]

On peut aussi écrire :

\[
25\,\%
=
\dfrac{1}{4}.
\]

De même :

\[
50\,\%
=
\dfrac{50}{100}
=
\dfrac{1}{2}.
\]

Et :

\[
10\,\%
=
\dfrac{10}{100}
=
\dfrac{1}{10}.
\]

Le symbole $\%$ signifie donc « pour cent ».

\section{25. Calculer une proportion}

Supposons qu’un groupe compte $40$ élèves et que $10$ d’entre eux choisissent une activité.

La proportion est :

\[
\dfrac{10}{40}.
\]

On obtient :

\[
\dfrac{10}{40}
=
\dfrac{1}{4}.
\]

Comme :

\[
\dfrac{1}{4}
=
25\,\%,
\]

la proportion est :

\[
\boxed{25\,\%}.
\]

Dans les cas simples, on utilise des fractions usuelles ou des écritures équivalentes à une fraction de dénominateur $100$.

\section{26. Appliquer un pourcentage}

Calculer $25\,\%$ de $80$ revient à calculer :

\[
\dfrac{25}{100}\times80.
\]

Comme :

\[
\dfrac{25}{100}
=
\dfrac{1}{4},
\]

on obtient :

\[
80\div4=20.
\]

Donc :

\[
\boxed{
25\,\%\text{ de }80=20
}.
\]

Le pourcentage agit ici comme une fraction opérateur.

\section{27. Exemple développé 3 — comparer deux proportions}

Dans un premier groupe, $8$ élèves sur $20$ réussissent une tâche.

La proportion vaut :

\[
\dfrac{8}{20}
=
\dfrac{40}{100}
=
40\,\%.
\]

Dans un second groupe, $9$ élèves sur $30$ réussissent.

La proportion vaut :

\[
\dfrac{9}{30}
=
\dfrac{3}{10}
=
30\,\%.
\]

Même si $9>8$, la proportion de réussite est plus faible dans le second groupe.

Cette situation montre pourquoi il faut comparer des \textbf{proportions} et pas seulement des effectifs.

\section{28. Automatismes à maîtriser}

Certaines relations doivent devenir rapides.

\subsection{Fractions usuelles}

\[
\dfrac{1}{2}=0{,}5,
\]

\[
\dfrac{1}{4}=0{,}25,
\]

\[
\dfrac{3}{4}=0{,}75,
\]

\[
\dfrac{3}{2}=1{,}5,
\]

\[
\dfrac{5}{2}=2{,}5.
\]

\subsection{Relations entre fractions usuelles}

\[
\dfrac{1}{2}+\dfrac{1}{2}=1,
\]

\[
\dfrac{1}{2}+\dfrac{1}{4}=\dfrac{3}{4},
\]

\[
1-\dfrac{1}{4}=\dfrac{3}{4},
\]

\[
\dfrac{3}{4}-\dfrac{1}{4}=\dfrac{1}{2}.
\]

\section{29. Méthode — choisir une représentation utile}

Les fractions peuvent être représentées de plusieurs façons :

\begin{itemize}
\item bande ou rectangle partagé ;
\item disque partagé ;
\item demi-droite graduée ;
\item tableau de proportion ;
\item écriture fractionnaire ;
\item écriture décimale lorsque c’est possible ;
\item pourcentage dans les cas adaptés.
\end{itemize}

La représentation choisie dépend de la question.

Pour comprendre un partage, une bande peut être très utile.

Pour comparer des nombres, une demi-droite graduée peut être plus efficace.

Pour calculer une proportion, une fraction ou un tableau peut être préférable.

\section{30. Erreurs fréquentes}

\subsection{Penser qu’une fraction est toujours inférieure à $1$}

C’est faux :

\[
\dfrac{7}{4}>1.
\]

\subsection{Confondre numérateur et dénominateur}

Dans :

\[
\dfrac{3}{5},
\]

le dénominateur $5$ indique le nombre de parts égales d’une unité ; le numérateur $3$ indique combien de parts sont prises.

\subsection{Additionner les dénominateurs}

Il est faux d’écrire :

\[
\dfrac{1}{5}+\dfrac{2}{5}
=
\dfrac{3}{10}.
\]

On doit écrire :

\[
\dfrac{1}{5}+\dfrac{2}{5}
=
\dfrac{3}{5}.
\]

\subsection{Confondre fraction et quotient approché}

On peut avoir :

\[
\dfrac{2}{3}\approx0{,}67,
\]

mais :

\[
\dfrac{2}{3}\neq0{,}67.
\]

La fraction est exacte ; l’écriture décimale est ici seulement approchée.

\subsection{Appliquer un pourcentage sans identifier le tout}

Pour calculer une proportion ou un pourcentage, il faut savoir ce qui constitue le \textbf{tout}.

\subsection{Utiliser le produit en croix}

Le produit en croix n’est pas la méthode attendue dans ce chapitre de 6e. On privilégie le sens des fractions, les fractions égales, le retour à l’unité et les relations simples.

\section{31. À retenir}

Une fraction est un \textbf{nombre}.

Elle peut représenter un quotient :

\[
\boxed{
\dfrac{a}{b}=a\div b
}
\]

avec $b\neq0$.

Elle vérifie :

\[
\boxed{
b\times\dfrac{a}{b}=a
}.
\]

Elle peut aussi agir comme opérateur :

\[
\boxed{
\dfrac{a}{b}\times n
}.
\]

Des fractions différentes peuvent représenter le même nombre :

\[
\boxed{
\dfrac{a}{b}
=
\dfrac{a\times k}{b\times k}
}
\]

pour $k\neq0$.

Enfin, un pourcentage est une proportion exprimée pour $100$ :

\[
\boxed{
p\,\%=\dfrac{p}{100}
}.
\]

Le cœur du chapitre consiste à comprendre \textbf{ce que signifie la fraction dans chaque situation}, puis à choisir la représentation et le calcul adaptés.
